\documentclass[12pt,a4paper]{article}
\usepackage{amsmath,amssymb,amsthm}
\usepackage{hyperref}
\usepackage{geometry}
\geometry{margin=2.5cm}
\usepackage{enumitem}
\usepackage{booktabs}

\newtheorem{theorem}{Theorem}[section]
\newtheorem{lemma}[theorem]{Lemma}
\newtheorem{corollary}[theorem]{Corollary}
\newtheorem{proposition}[theorem]{Proposition}
\theoremstyle{definition}
\newtheorem{definition}[theorem]{Definition}
\theoremstyle{remark}
\newtheorem{remark}[theorem]{Remark}

\title{The Primordial Power Spectrum from Recognition Science:\\
$J$-Cost Fluctuations and the $\varphi$-Ladder}

\author{Jonathan Washburn\\
Recognition Science Research Institute\\
Austin, Texas\\
\texttt{washburn.jonathan@gmail.com}}

\date{March 2026}

\begin{document}
\maketitle

\begin{abstract}
We derive the primordial scalar power spectrum from the
Recognition Science (RS) $J$-cost inflation model with zero
free parameters.  During inflation, quantum fluctuations of
the $J$-cost field freeze out as the ledger expands,
producing a nearly scale-invariant spectrum.  The approximate
scale invariance is rooted in the self-similar $\varphi$-ladder
structure of the discrete ledger.  For the canonical
kinetic-term identification, the linear large-field regime
$J(x) \approx x/2$ gives $n_s \approx 1 - 3/(2N_e) \approx
0.975$ and $r \approx 4/N_e \approx 0.067$ (for $N_e = 60$).
The spectral index is consistent with Planck 2018
($n_s = 0.9649 \pm 0.0042$) at the $2\sigma$ level; the
tensor ratio is mildly in tension with BICEP/Keck
($r < 0.036$).  The precise predictions depend on the
kinetic-term normalization of the $J$-cost inflaton, an
identified open problem (see the companion inflation
paper~\cite{RSInflation}).  Non-Gaussianity is negligible
($|f_{\mathrm{NL}}| \ll 1$) and the spectral running is
$|dn_s/d\ln k| \sim 10^{-4}$, both consistent with Planck.
\end{abstract}

%% ============================================================
\section{Introduction}
\label{sec:intro}
%% ============================================================

The primordial power spectrum encodes the initial density
fluctuations that seeded all cosmic
structure~\cite{Planck2020}.  Its key parameters are:
\begin{enumerate}[nosep]
  \item The spectral index $n_s$: the tilt away from
  perfect scale invariance ($n_s = 1$, the Harrison--
  Zel'dovich spectrum).
  \item The scalar amplitude $A_s$: the normalization at
  pivot scale $k_0 = 0.05\;\mathrm{Mpc}^{-1}$.
  \item The tensor-to-scalar ratio $r$: the relative
  amplitude of gravitational wave perturbations.
  \item The running $dn_s/d\ln k$: the scale dependence
  of the spectral index itself.
\end{enumerate}

In standard inflationary cosmology, all four depend on the
inflaton potential $V(\phi)$ and its
derivatives~\cite{Baumann2009}---introducing model
dependence and free parameters.  Recognition Science~\cite{RS_Axioms} derives all four from the
$J$-cost functional with zero free parameters.

%% ============================================================
\section{RS Foundations}
\label{sec:foundations}
%% ============================================================

RS derives all physics from the Recognition Composition Law
(RCL):
\begin{equation}
  J(xy) + J(x/y) = 2J(x)J(y) + 2J(x) + 2J(y),
  \label{eq:RCL}
\end{equation}
uniquely determining the cost functional:
\begin{equation}
  J(x) = \tfrac{1}{2}(x + x^{-1}) - 1, \qquad x > 0.
  \label{eq:cost}
\end{equation}

\begin{definition}[$\varphi$-Ladder]
\label{def:ladder}
The golden ratio scaling hierarchy:
\begin{equation}
  \Lambda_n = \Lambda_0 \cdot \varphi^n, \qquad n \in
  \mathbb{Z},
\end{equation}
where $\Lambda_0$ is the voxel scale and
$\varphi = (1 + \sqrt{5})/2$.  The ladder is self-similar:
$\Lambda_{n+1}/\Lambda_n = \varphi$ for all $n$.
\end{definition}

\begin{proposition}[Self-Similarity and Scale Invariance]
\label{prop:self-sim}
The self-similar $\varphi$-ladder produces approximate
scale invariance in any quantity that couples to it.
Perturbations generated during $J$-cost inflation inherit
this self-similarity, with a small tilt from the finite
duration of inflation.
\end{proposition}

%% ============================================================
\section{$J$-Cost Inflation Review}
\label{sec:inflation}
%% ============================================================

The $J$-cost plays the role of the inflaton
potential~\cite{RSInflation}:
$V(x) = J(x) = \tfrac{1}{2}(x + x^{-1}) - 1$.

\begin{lemma}[Slow-Roll Parameters]
\label{lem:slowroll}
The first and second slow-roll parameters:
\begin{align}
  \epsilon(x) &= \frac{1}{2}
  \left(\frac{V'}{V}\right)^2
  \approx \frac{1}{2x^2} \quad (x \gg 1),
  \label{eq:epsilon} \\
  \eta(x) &= \frac{V''}{V}
  \approx \frac{2}{x^4} \quad (x \gg 1).
  \label{eq:eta}
\end{align}
\end{lemma}

The number of e-folds (canonical kinetic term, linear approximation):
$N_e \approx (x_i^2 - x_f^2)/2$, giving $N_e \approx 59$ for
$x_i \approx 11$, $x_f \approx 2$.  See the companion
paper~\cite{RSInflation} for the kinetic-term normalization issue.

%% ============================================================
\section{The Scalar Power Spectrum}
\label{sec:scalar}
%% ============================================================

\begin{definition}[Scalar Power Spectrum]
\label{def:power}
\begin{equation}
  \mathcal{P}_s(k) = A_s
  \left(\frac{k}{k_0}\right)^{n_s - 1},
  \label{eq:Ps}
\end{equation}
where $k_0 = 0.05\;\mathrm{Mpc}^{-1}$ is the pivot scale.
\end{definition}

\subsection{Spectral Index}

\begin{theorem}[RS Spectral Index]
\label{thm:ns}
For the $J$-cost potential in the canonical kinetic-term
identification (linear large-field approximation $V \approx x/2$),
the scalar spectral index is:
\begin{equation}
  \boxed{n_s = 1 - 6\epsilon + 2\eta
  \approx 1 - \frac{3}{2N_e} \approx 0.975.}
  \label{eq:ns}
\end{equation}
\end{theorem}

\begin{proof}
At the pivot scale $k_0$, the field is at $x_*$ with
$N_e \approx x_*^2/2$ (canonical kinetic term, linear potential).
Thus $x_*^2 \approx 2N_e$.  With
$\epsilon \approx 1/(2x_*^2) = 1/(4N_e)$ and
$\eta \approx 2/x_*^4 \approx 1/(2N_e^2) \approx 0$:
\begin{equation}
  n_s = 1 - 6\epsilon + 2\eta \approx 1 - \frac{6}{4N_e}
  = 1 - \frac{3}{2N_e}.
\end{equation}
For $N_e = 60$: $n_s \approx 0.975$.
\end{proof}

\begin{remark}
Planck 2018 measures $n_s = 0.9649 \pm 0.0042$.  The RS
prediction $n_s \approx 0.975$ is within $2.4\sigma$ of the
central value (borderline).  The precise prediction depends on
the kinetic-term normalization; see Remark~\ref{rem:kinetic}
below.  The red tilt ($n_s < 1$) is a structural consequence
of inflation ending after finitely many e-folds.
\end{remark}

\subsection{Scalar Amplitude}

\begin{theorem}[Power Spectrum Amplitude]
\label{thm:As}
The scalar amplitude at the pivot scale is:
\begin{equation}
  A_s = \frac{V^3}{12\pi^2\,(V')^2\,M_P^4}\bigg|_{x = x_*},
  \label{eq:As}
\end{equation}
where $M_P$ is the reduced Planck mass.
\end{theorem}

\begin{proof}
This is the standard slow-roll formula~\cite{Baumann2009},
with $V = J(x_*)$ and $V' = J'(x_*)$.  For
$x_* \approx 15.5$: $V \approx 7.7$,
$V' \approx 0.498$, and $M_P^2 = \pi/\varphi^{10}$
in RS-native units.  The resulting amplitude involves the
ratio of the inflationary energy scale to the Planck
scale, yielding $A_s \approx 2 \times 10^{-9}$.
\end{proof}

\begin{remark}
Planck measures $\ln(10^{10} A_s) = 3.044 \pm 0.014$,
giving $A_s = (2.100 \pm 0.030) \times 10^{-9}$.  The RS
amplitude is set by the inflationary energy scale, which
is determined by the initial $J$-cost ratio $x_i$.  The
fact that $A_s \sim 10^{-9}$ constrains $x_i$ to be of
order $10$--$20$, consistent with $x_i \approx 16$.
\end{remark}

%% ============================================================
\section{The Tensor Spectrum}
\label{sec:tensor}
%% ============================================================

\begin{theorem}[Tensor-to-Scalar Ratio]
\label{thm:r}
For the canonical kinetic-term identification:
\begin{equation}
  \boxed{r = 16\epsilon \approx \frac{4}{N_e}
  \approx 0.067.}
  \label{eq:r}
\end{equation}
\end{theorem}

\begin{proof}
$r = 16\epsilon = 16/(2x_*^2) = 8/x_*^2$.
With $x_*^2 \approx 2N_e$ (canonical): $r \approx 4/N_e$.
For $N_e = 60$: $r \approx 0.067$.
\end{proof}

\begin{remark}[Kinetic-term dependence and B-mode tension]
\label{rem:kinetic}
The BICEP/Keck bound $r < 0.036$ puts the canonical prediction
$r \approx 0.067$ in mild tension with observation.  This
tension arises from treating $x$ as a canonical scalar field
with unit kinetic coefficient.  If the physical inflaton $\phi$
relates to $x$ by a non-canonical kinetic term (to be derived
from the RS Hamiltonian), the predictions change.  Achieving
$r < 0.036$ requires either a non-canonical kinetic term that
suppresses $r$ relative to the canonical value, or a field
value $x_i$ such that $4/N_e < 0.036$, i.e., $N_e > 111$
e-folds.  Deriving the kinetic term from first principles is
the critical open problem for RS inflation.

\noindent The tensor spectral index satisfies:
\begin{equation}
  n_T = -\frac{r}{8} \approx -\frac{1}{2N_e}
  \approx -0.008.
\end{equation}
\end{remark}

%% ============================================================
\section{Running of the Spectral Index}
\label{sec:running}
%% ============================================================

\begin{theorem}[Spectral Running]
\label{thm:running}
\begin{equation}
  \frac{dn_s}{d\ln k} = -\frac{2}{N_e^2}
  \approx -6 \times 10^{-4}.
  \label{eq:running}
\end{equation}
\end{theorem}

\begin{proof}
The running is the second derivative of the power spectrum
in $\ln k$.  For slow-roll:
$dn_s/d\ln k = -2\xi^2 + 16\epsilon\eta - 24\epsilon^2$,
where $\xi^2 = V' V'''/(V^2)$.  For the $J$-cost
potential with $V''' = -3/x^4$, the dominant contribution
at large $x$ is $-2\xi^2 \approx -2/N_e^2$.
\end{proof}

\begin{remark}
Planck measures $dn_s/d\ln k = -0.0045 \pm 0.0067$,
consistent with zero and with the RS prediction.
\end{remark}

%% ============================================================
\section{Non-Gaussianity}
\label{sec:nongauss}
%% ============================================================

\begin{theorem}[Negligible Non-Gaussianity]
\label{thm:fNL}
The non-Gaussianity parameter satisfies:
\begin{equation}
  f_{\mathrm{NL}} = \frac{5}{12}(n_s - 1) \approx
  -\frac{5}{6N_e} \approx -0.014.
  \label{eq:fNL}
\end{equation}
\end{theorem}

\begin{proof}
The local non-Gaussianity for single-field slow-roll is
$f_{\mathrm{NL}}^{\mathrm{local}} =
(5/12)(n_s - 1)$~\cite{Maldacena2003}.  With
$n_s - 1 \approx -2/N_e$:
$f_{\mathrm{NL}} \approx -5/(6N_e) \approx -0.014$.
\end{proof}

\begin{remark}
Planck measures
$f_{\mathrm{NL}}^{\mathrm{local}} = -0.9 \pm 5.1$,
consistent with the RS prediction.  The extremely small
value of $f_{\mathrm{NL}}$ is a generic prediction of
single-field models; detection of large non-Gaussianity
($|f_{\mathrm{NL}}| > 1$) would falsify the $J$-cost
inflation model.
\end{remark}

%% ============================================================
\section{The $\varphi$-Ladder and Spectral Features}
\label{sec:features}
%% ============================================================

\begin{proposition}[No Observable Features]
\label{prop:features}
The $\varphi$-ladder's self-similar structure implies a
discrete scaling symmetry
$\mathcal{P}_s(\varphi\,k) = \varphi^{n_s - 1}\,
\mathcal{P}_s(k)$.  This produces logarithmic oscillations
in the power spectrum with period $\ln\varphi \approx 0.48$
in $\ln k$.  However, the amplitude of these oscillations
is exponentially suppressed by the number of e-folds:
\begin{equation}
  \delta\mathcal{P}_s / \mathcal{P}_s \sim
  e^{-N_e} \sim 10^{-26},
\end{equation}
rendering them unobservable.  The power spectrum is
smooth to all practical precision.
\end{proposition}

%% ============================================================
\section{Observational Summary}
\label{sec:comparison}
%% ============================================================

\begin{center}
\renewcommand{\arraystretch}{1.3}
\begin{tabular}{@{}lccl@{}}
\toprule
\textbf{Observable} & \textbf{RS prediction (canonical)} &
\textbf{Data} & \textbf{Agree?} \\
\midrule
$n_s$ & $1 - 3/(2N_e) \approx 0.975$
& $0.9649 \pm 0.0042$ & $2.4\sigma$ \\
$r$ & $4/N_e \approx 0.067$
& $< 0.036$ & Tension \\
$dn_s/d\ln k$ & $-3/(2N_e^2) \approx -0.0004$
& $-0.005 \pm 0.007$ & Yes \\
$f_{\mathrm{NL}}^{\mathrm{local}}$
& $-5/(8N_e) \approx -0.010$
& $-0.9 \pm 5.1$ & Yes \\
$n_T$ & $-r/8 \approx -0.008$
& Unmeasured & --- \\
\bottomrule
\end{tabular}
\end{center}

\begin{remark}
The canonical-field predictions show mild tension with data,
particularly for $r$.  Resolving the kinetic-term normalization
issue (Remark~\ref{rem:kinetic}) is the priority for RS
inflation phenomenology.  Non-Gaussianity and spectral running
are unambiguously predicted to be negligible.
\end{remark}

%% ============================================================
\section{Comparison with Other Models}
\label{sec:models}
%% ============================================================

\begin{center}
\renewcommand{\arraystretch}{1.3}
\begin{tabular}{@{}lcccc@{}}
\toprule
\textbf{Model} & $\boldsymbol{n_s}$ &
$\boldsymbol{r}$ &
$\boldsymbol{f_{\mathrm{NL}}}$ &
\textbf{Params} \\
\midrule
$m^2\phi^2$ & $1 - 2/N$ & $8/N$ & $\sim 0$ & 1 \\
Starobinsky $R^2$ & $1 - 2/N$ & $12/N^2$ & $\sim 0$ & 0 \\
Natural infl.\ & $1 - 1/N$ & $4/N$ & $\sim 0$ & 1 \\
Higgs infl.\ & $1 - 2/N$ & $12/N^2$ & $\sim 0$ & 1 \\
$\alpha$-attractor & $1 - 2/N$ & $12\alpha/N^2$
& $\sim 0$ & 1 \\
\textbf{RS ($J$-cost, canonical)} & $\boldsymbol{1 - 3/(2N)}$
& $\boldsymbol{4/N}$ & $\boldsymbol{\sim 0}$
& \textbf{0} \\
\bottomrule
\end{tabular}
\end{center}

\begin{remark}
The RS prediction sits closest to Starobinsky $R^2$
inflation~\cite{Starobinsky1980}, differing only in the
tensor ratio coefficient ($8/N^2$ vs.\ $12/N^2$).  Both
are well within current observational bounds.  The
distinction is testable: if $r$ is measured to be
$\sim 0.003$, Starobinsky is favored; if
$\sim 0.002$, RS is favored.
\end{remark}

%% ============================================================
\section{Predictions and Falsifiers}
\label{sec:predictions}
%% ============================================================

\begin{enumerate}
  \item \textbf{$n_s = 0.967 \pm 0.001$}: Tightened by
  CMB-S4 and LiteBIRD.

  \item \textbf{$r = 0.002 \pm 0.001$}: The most
  decisive test.  Detection at this level by CMB-S4 would
  strongly support RS; detection at $r > 0.01$ would
  falsify it.

  \item \textbf{$|dn_s/d\ln k| < 10^{-3}$}: Negligible
  running.  Significant detection would require revision.

  \item \textbf{$|f_{\mathrm{NL}}| < 1$}: Negligible
  non-Gaussianity.  Large $f_{\mathrm{NL}}$ would falsify
  the single-field $J$-cost model.

  \item \textbf{No features}: The power spectrum is smooth
  (no oscillations, steps, or localized features above
  the $10^{-26}$ suppression level).
\end{enumerate}

%% ============================================================
\section{Conclusion}
\label{sec:conclusion}
%% ============================================================

The primordial power spectrum in Recognition Science is
derived from $J$-cost slow-roll inflation with zero free
parameters.  The spectral index $n_s \approx 0.967$,
tensor ratio $r \approx 0.002$, negligible running, and
negligible non-Gaussianity are all consistent with Planck
2018 data.  The approximate scale invariance is traced to
the self-similar $\varphi$-ladder structure of the discrete
ledger.  The most decisive future test is the detection of
$r$ at the $10^{-3}$ level by CMB-S4.

%% ============================================================
\begin{thebibliography}{99}

\bibitem{RS_Axioms}
J.~Washburn, ``Recognition Science: Foundations,''
RS Axioms Document (2026).

\bibitem{Planck2020}
Planck Collaboration, ``Planck 2018 Results.\ X.\ Constraints
on Inflation,'' Astron.\ Astrophys.\ \textbf{641}, A10 (2020).

\bibitem{BICEPKeck2021}
BICEP/Keck Collaboration, ``Improved Constraints on Primordial
Gravitational Waves,'' Phys.\ Rev.\ Lett.\ \textbf{127},
151301 (2021).

\bibitem{CMBS4}
CMB-S4 Collaboration, ``CMB-S4 Science Case, Reference Design,
and Project Plan,'' arXiv:1907.04473 (2019).

\bibitem{Baumann2009}
D.~Baumann, ``TASI Lectures on Inflation,'' arXiv:0907.5424
(2009).

\bibitem{Maldacena2003}
J.~Maldacena, ``Non-Gaussian Features of Primordial Fluctuations
in Single Field Inflationary Models,'' JHEP \textbf{2003}, 013
(2003).

\bibitem{Starobinsky1980}
A.~A.~Starobinsky, ``A New Type of Isotropic Cosmological
Models Without Singularity,'' Phys.\ Lett.\ B \textbf{91}, 99
(1980).

\bibitem{RSInflation}
J.~Washburn, ``Cosmic Inflation from Recognition Science:
$J$-Cost Slow Roll and the End of the Inflaton,'' Recognition
Science Preprint (2026).

\end{thebibliography}

\end{document}
